\begin{lemma}
For every fixed \(\delta>0\) and \(\varepsilon>0\), there is a threshold
\(n_0\) such that the following holds for every \(n\ge n_0\).  Let
\[
m=\noderef{TwoBiteNaturalReducedVertexCount}(n)
=\left\lfloor\frac{n}{\log^2 n}\right\rfloor,
\qquad
k=\noderef{TwoBiteNaturalIndependenceScale}(1+\varepsilon,n)
  =\left\lceil(1+\varepsilon)\sqrt{n\log n}\right\rceil .
\]
For integers \(\ell_R,\ell_B\) with \(\ell_R\ell_B\ge k\), put
\[
x_R=\frac{\ell_R}{k},\qquad x_B=\frac{\ell_B}{k}.
\]
If \(2\delta k < (2-x_R-x_B)k\), then
\[
\frac{\binom m{\ell_R}\binom m{\ell_B}\binom{\ell_R\ell_B}{k}}{\binom{m^2}{k}}
\le
\exp\left(\left(\delta-\frac{2-x_R-x_B}{2}\right)k\log n\right).
\]
\end{lemma}
\begin{proof}
Write \(L=\log n\) and
\[
s=(2-x_R-x_B)k.
\]
The hypothesis says \(s>2\delta k\).  In particular \(k>0\), \(s>0\), and
\[
x_R+x_B<2,\qquad 0<s\le 2k,
\]
because \(x_R,x_B\ge0\).

Increase \(n_0\), if necessary, so that the conclusions of
\(\noderef{ProjectionSizeAnalyticParameterBounds}\) hold and also
\[
L\ge 1,\qquad
2\left(\frac52\log L+\log(4(1+\varepsilon))\right)+3\le \delta L.
\]
This is possible because \(\log L=o(L)\).  The parameter bound gives
\[
m\ge \frac{n}{2L^2},\qquad
k\le 2(1+\varepsilon)\sqrt{nL},\qquad
1\le k\le n,\qquad n\le m^2.
\]
Thus \(\binom{m^2}{k}>0\), so the denominator is positive.  If
\(\ell_R>m\) or \(\ell_B>m\), then one of
\(\binom m{\ell_R}\), \(\binom m{\ell_B}\) is zero, and the desired
inequality is immediate.  Hence assume \(\ell_R,\ell_B\le m\).  Since
\(\ell_R\ell_B\ge k\ge1\), both \(\ell_R\) and \(\ell_B\) are positive, so
all logarithms below are legitimate.

We use the elementary estimates
\[
\binom ab\le \left(\frac{ea}{b}\right)^b
\quad(1\le b\le a),
\qquad
\binom Nk\ge \left(\frac Nk\right)^k
\quad(1\le k\le N).
\]
Applying these to the three numerator factors and the denominator gives
\[
\frac{\binom m{\ell_R}\binom m{\ell_B}\binom{\ell_R\ell_B}{k}}
{\binom{m^2}{k}}
\le
\frac{
\left(\frac{em}{\ell_R}\right)^{\ell_R}
\left(\frac{em}{\ell_B}\right)^{\ell_B}
\left(\frac{e\ell_R\ell_B}{k}\right)^k}
\left(\frac{m^2}{k}\right)^k}.
\]
Taking logarithms and using
\(\ell_R=kx_R\), \(\ell_B=kx_B\), we obtain
\[
\begin{aligned}
\log R
&\le
kx_R\left(1+\log\frac{m}{k}-\log x_R\right)
+kx_B\left(1+\log\frac{m}{k}-\log x_B\right)  \\
&\quad
+k\left(1+\log k+\log x_R+\log x_B\right)
-k\left(2\log m-\log k\right),
\end{aligned}
\]
where \(R\) denotes the binomial ratio on the left.  Since
\(2\log k-2\log m=-2\log(m/k)\), this simplifies to
\[
\log R
\le
k(1+x_R+x_B)
-s\log\frac{m}{k}
+k(1-x_R)\log x_R
+k(1-x_B)\log x_B.
\]
For every \(x>0\), \((1-x)\log x\le0\): it is negative for \(x<1\), zero at
\(x=1\), and negative for \(x>1\).  Therefore
\[
\log R\le k(1+x_R+x_B)-s\log\frac{m}{k}
=(3k-s)-s\log\frac{m}{k}.
\]

It remains to lower-bound \(\log(m/k)\).  From the parameter estimates,
\[
\frac{m}{k}
\ge
\frac{n/(2L^2)}{2(1+\varepsilon)\sqrt{nL}}
=
\frac{\sqrt n}{4(1+\varepsilon)L^{5/2}}.
\]
Thus
\[
\log\frac{m}{k}
\ge
\frac12 L-\frac52\log L-\log(4(1+\varepsilon)).
\]
Let
\[
E=\frac52\log L+\log(4(1+\varepsilon)).
\]
Since \(L\ge1\) and \(\varepsilon>0\), we have \(E\ge0\).
Then
\[
\log R
\le
3k-s-s\left(\frac12 L-E\right)
=-\frac{s}{2}L+(sE+3k-s).
\]
Because \(0<s\le2k\), the error term satisfies
\[
sE+3k-s\le 2kE+3k\le \delta kL
\]
by the choice of \(n_0\).  Hence
\[
\log R\le -\frac{s}{2}L+\delta kL
=\left(\delta-\frac{2-x_R-x_B}{2}\right)k\log n.
\]
Exponentiating both sides, using monotonicity of \(\exp\), gives exactly
\[
R\le
\exp\left(\left(\delta-\frac{2-x_R-x_B}{2}\right)k\log n\right).
\]
\end{proof}
